\section{From the certificate to the exponent bound}\label{sec:conclusion}

We now pass from the strict comparison to the exponent bound.
This last step uses the finite construction class defining the attainable
value, rather than a limiting tensor.

\begin{proof}[Proof of the square assertion in \cref{thm:main}]
Set \(\beta=1129/500\).
By \cref{thm:certificate}, at the center \(\alpha_*\) of the certified chart,
\[
F_\beta(\alpha_*)\ge 1.098619383270680538
>\log3+7.094\times10^{-6}.
\]
By the definition of \(F_\beta\) as the closure of scores of admissible
finite constructions, there is a finite construction whose score is
strictly greater than \(\log3\).
Its law may approximate \(\alpha_*\); exact equality of the law is not
needed in the rank inequality.

Let its six banks contain \(N\) original occurrences each and put \(n=6N\).
Let \(L\) be its number of nonzero direct labels,
\(R\) its charged ancillary border-rank bound, and
\(v=abc\) its common matrix multiplication volume.
The construction gives a finite degeneration
\[
T^{\otimes n}\otimes\mathcal A
\unrhd\bigoplus_{\ell=1}^L\langle a,b,c\rangle,
\qquad \brank(\mathcal A)\le R.
\]
All parameters here---including group sizes, control horizons, rational
approximants and exact bank sizes---have already been chosen.
The score inequality reads
\[
\frac{\log L-\log R+(\beta/3)\log v}{n}>\log3,
\]
or \(Lv^{\beta/3}>3^nR\).
On the other hand, \cref{thm:asi} and \(\brank(T)=3\) give
\(Lv^{\omega/3}\le3^nR\).
Since \(v\ge1\), the inequality \(\omega\ge\beta\) is impossible.
Therefore \(\omega<1129/500\).
\end{proof}

For the rational exponent in this paper, the sufficient finite inequality
has the equivalent integer form
\[
L^{1500}v^{1129}>(3^nR)^{1500}.
\]
Thus the strict comparison does not depend on taking an asymptotic sum
inequality at an infinite tensor or discarding an ancillary charge.
The argument establishes an arithmetic exponent.  It does not estimate
the size of a useful finite implementation or its numerical stability.
